\section{\texorpdfstring{$\Ham(M,\omega)$}{Ham(M, omega)} 上的霍费尔测地线}
\stfhIndex{I163}{huo^fei^er^ce^de^xian}{霍费尔测地线（Hofer’s geodesic）}
\label{sec:hofer-geodesics}

现在设 $\phi,\psi\in\Ham(M,\omega)$，并考虑霍费尔距离
$d(\phi,\psi)$。利用霍费尔距离的不变性，不失一般性，可以假设
$\psi=\operatorname{id}$。从 $\Ham(M,\omega)$ 的霍费尔芬斯勒几何观点看，一个自然的问题是下确界
\[
 \|\phi\|=\inf\{\|H\|\mid H\mapsto\phi\}
\]
是否能达到；也就是说，是否存在哈密顿量 $H$，使得
$\|\phi\|=\|H\|$。

\begin{definition}\label{def:hamiltonian-equivalence}
\stfhIndex{I015}{00-symbol-015}{$F\sim G$}
\stfhIndex{I016}{00-symbol-016}{$[F]$}
\stfhIndex{I049}{00-symbol-049}{$\widetilde\phi_F$}
如果与两个哈密顿量 $F$ 和 $G$ 相伴的哈密顿路径 $\phi_F$ 与
$\phi_G$ 相对于边界 $t=0,1$ 道路同伦，就称 $F$ 与 $G$ 等价，并记
$F\sim G$。等价地，如果两个哈密顿量由单参数哈密顿量族
$\{F^s\}_{0\leq s\leq1}$ 连接，并且对所有 $s\in[0,1]$ 都有
$F^s\mapsto\phi$，则 $F\sim G$。以 $[F]$ 或
$\widetilde\phi_F$ 表示 $F$ 的等价类。再以
$\widetilde{\Ham}(M,\omega)$ 表示 $\Ham(M,\omega)$ 的这类等价类的集合。
\end{definition}

\begin{definition}\label{def:hofer-length-minimizing}
\stfhSubIndex{I164}{huo^fei^er^ce^de^xian}{霍费尔测地线（Hofer’s geodesic）}{chang^du^ji^xiao}{长度极小（length minimizing）}
\stfhSubIndex{I165}{huo^fei^er^ce^de^xian}{霍费尔测地线（Hofer’s geodesic）}{zai^qi^tong^lun^lei^zhong^chang^du^ji^xiao}{在其同伦类中长度极小（length minimizing in its homotopy class）}
设 $\phi_H:t\in[0,1]\mapsto\phi_H^t$ 是一条满足
$\phi_H^1=\phi\in\Ham(M,\omega)$ 的哈密顿路径。
\begin{enumerate}[label=(\arabic*)]
 \item 如果对任意满足
 \[
  \phi_{H'}^0=\phi_H^0,
  \qquad
  \phi_{H'}^1=\phi_H^1
 \]
 的哈密顿路径 $\phi_{H'}$，都有
 $\operatorname{leng}(\phi_{H'})\geq\operatorname{leng}(\phi_H)$，或者等价地，$\|H'\|\geq\|H\|$，就称 $\phi_H$（或 $H$）是\emph{长度极小的}。
 \item 如果对另外还满足 $H\sim H'$ 的 $H'$，同一结论成立，就称
 $\phi_H$ 在其同伦类中是\emph{长度极小的}。
\end{enumerate}
\end{definition}

霍费尔 (H93) 证明：只要相应的哈密顿方程\emph{没有周期小于或等于 $1$ 的非常值周期轨道}，$\mathbb C^n$ 上任意紧支\emph{自治}哈密顿量的路径都是长度极小的。

可以与黎曼几何类似地定义 $\Ham(M,\omega)$ 上的\emph{测地线}概念。然而，由于霍费尔长度函数并不光滑，给出测地线的精确定义并非完全平凡。
(Po01) 给出了 $\Ham(M,\omega)$ 上测地线的精确变分定义；关于测地线的一些一般性讨论，我们也请读者参阅该文。

\stfhIndex{I062}{bi^ya^li^bo^er^jie^luo^wei^qi}{比亚利--波尔捷罗维奇（Bialy–Polterovich）}
这里只引用米沙·比亚利（Misha Bialy）与波尔捷罗维奇 (BP94) 用相应哈密顿量的动力学性质给出的一个测地线刻画。

\begin{definition}\label{def:quasi-autonomous-hamiltonian}
如果存在两点 $x_-,x_+\in M$，使得对所有 $t\in[0,1]$ 都有
\[
 H(t,x_-)=\min_xH(t,x),
 \qquad
 H(t,x_+)=\max_xH(t,x),
\]
就称哈密顿量 $H$ 是\emph{准自治的}。
\stfhSubIndex{I156}{ha^mi^dun}{哈密顿（Hamiltonian）}{zhun^zi^zhi}{准自治（quasi-autonomous）}
\end{definition}

在 (BP94) 与 (LM95b) 中，比亚利与波尔捷罗维奇以及拉隆德与麦克达夫证明了：任意长度极小的（相应地，局部长度极小的）哈密顿路径都由\emph{准自治}的（相应地，\emph{局部准自治}的）哈密顿路径生成。

现在回忆伊利亚·乌斯季洛夫斯基（Ilya Ustilovsky）、拉隆德与麦克达夫关于测地线稳定性给出的必要条件。
\stfhSubIndex{I167}{huo^fei^er^ce^de^xian}{霍费尔测地线（Hofer’s geodesic）}{wu^si^ji^luo^fu^si^ji^la^long^de^mai^ke^da^fu^bi^yao^tiao^jian}{乌斯季洛夫斯基--拉隆德--麦克达夫必要条件（Ustilovsky–Lalonde–McDuff necessary condition）}

\begin{definition}\label{def:over-twisted-periodic-orbit}
设 $H:[0,1]\times M\to\mathbb R$ 是一个不必关于时间周期的哈密顿量，并设 $\phi_H^t$ 是其哈密顿流。
\begin{enumerate}[label=(\arabic*)]
 \item 如果 $\phi_H^T(p)=p$，就称点 $p\in M$ 是一个\emph{时间
 $T$-周期点}。如果道路
 $t\in[0,T]\mapsto\phi_H^t(p)$ 可缩，就称它是一个\emph{可缩时间
 $T$-周期轨道}。
 \item 当 $H$ 在 $t\in[0,T]$ 上有固定临界点 $p$ 时，如果
 $T_pM$ 上的线性化流
 $d\phi_H^t(p)$（$t\in[0,T]$）有一条周期小于或等于 $T$ 的闭轨道，就称作为时间 $T$-周期轨道的 $p$ 是\emph{过扭的}。
\end{enumerate}
\end{definition}

乌斯季洛夫斯基 (Ust96) 以及拉隆德与麦克达夫 (LM95b) 证明：对通有的
$\phi$（这里“通有”是指它的所有不动点都孤立），从单位元到 $\phi$ 的任意稳定测地线 $\phi_t$（$0\leq t\leq1$），都必须有至少两个欠扭的不动点。

关于相应自治路径长度极小性质的霍费尔定理，在 (En00)、(MSl01) 以及 (Oh02)--(Oh05b) 中得到推广；其附加假设是每个不动点处的线性化流都不过扭，即没有周期小于 $1$ 的闭轨道。下面的猜想由波尔捷罗维奇以 (Po01) 中猜想 12.6.D 的形式提出（也见 (LM95b) 与 (MSl01)），并由费利克斯·施伦克（Felix Schlenk） (Sch06) 与厄舍 (Ush10a) 以其最终形式证明。

\begin{theorem}[极小性猜想]\label{thm:minimality-conjecture}
\stfhSubIndex{I166}{huo^fei^er^ce^de^xian}{霍费尔测地线（Hofer’s geodesic）}{ji^xiao^xing^cai^xiang}{极小性猜想（minimality conjecture）}
\stfhIndex{I253}{pu^bu^bian^liang}{谱不变量（spectral invariants）}
任何没有周期小于或等于 $1$ 的可缩周期轨道的自治哈密顿路径，都在其相对于边界的道路同伦类中使霍费尔长度极小。
\end{theorem}

在本书第 4 部分，我们将通过\emph{谱不变量}的某种一般构造，把准自治路径的长度极小性质与相伴弗洛尔同调的某种同调性质联系起来，并用谱不变量给出长度极小性质的一个简单判据。利用这个判据，我们将依照 (Oh02)、(Oh05b) 与 (Ush10a) 所用方案，在第 4 部分给出极小性猜想的证明。

另一方面，$\Ham(M,\omega)$ 中并非每一对点都能由长度极小测地线连接，下面奥斯特罗弗 (Os03) 的定理说明了这一点。

\begin{theorem}[奥斯特罗弗]\label{thm:ostrover-no-minimizing-geodesic}
\stfhIndex{I222}{ao^si^te^luo^fu}{奥斯特罗弗（Ostrover）}
\stfhIndex{I192}{la^long^de}{拉隆德（Lalonde）}
\stfhIndex{I203}{mai^ke^da^fu}{麦克达夫（McDuff）}
设 $(M,\omega)$ 是满足 $\pi_2(M)=\{0\}$ 的辛流形。则存在元素
$\phi\in\Ham(M,\omega)$，它不能由一条长度极小测地线连接到单位元。
\end{theorem}

拉隆德与麦克达夫 (LM95b) 先前在 $S^2$ 上得到了这类结果。
